% ===========================================================================
% Category: Hidden Markov Model
% Generated from ML_Lab_Manual_Completed.md - edit the Markdown and re-run
% scripts/build_latex_manual.py instead of editing this file by hand.
% ===========================================================================

\chapter{Experiment 16: Hidden Markov Model
(HMM)}\label{experiment-16-hidden-markov-model-hmm}

\textbf{Category:} Probabilistic sequence modeling

\section{1. Objective}\label{objective}

Model a sequence with hidden states, transition probabilities and
emission distributions; decode the most likely hidden-state path with
Viterbi.

\section{2. Dataset Used}\label{dataset-used}

Real DAX stock index daily closes 1991-1998
(\passthrough{\lstinline!data/stock\_index.csv!}, R
\passthrough{\lstinline!EuStockMarkets!}), 1,859 trading days. The
observation sequence is the daily return; the close price is used only
for plotting. This is a realistic alternative to a toy discrete
sequence.

\section{3. Required Libraries}\label{required-libraries}

hmmlearn, numpy, pandas, matplotlib

\section{4. Brief Theory}\label{brief-theory}

An HMM assumes a hidden Markov state sequence z\_t (P(z\_t \textbar{}
z\_\{t−1\})) that emits observations x\_t \textbar{} z\_t = s\_k
\textasciitilde{} N(μ\_k, σ\_k²). The transition matrix A\_ij =
P(z\_\{t+1\}=j \textbar{} z\_t=i). Baum-Welch (EM) learns parameters
(forward-backward E-step, re-estimation M-step); Viterbi decoding finds
the most likely state path by dynamic programming; the forward algorithm
gives the log-likelihood.

\section{5. Procedure}\label{procedure}

\begin{enumerate}
\def\labelenumi{\arabic{enumi}.}
\tightlist
\item
  Load the DAX series; build the daily-return observation vector.
\item
  Fit
  \passthrough{\lstinline!GaussianHMM(n\_components=3, covariance\_type="full", n\_iter=200, random\_state=42)!}.
\item
  Check convergence and log-likelihood.
\item
  Decode the state path with Viterbi.
\item
  Sort states by volatility so they map to Bull/Sideways/Bear; reorder
  the transition matrix.
\item
  Plot the transition heatmap and the price colored by regime.
\end{enumerate}

\section{6. Reference Implementation}\label{reference-implementation}

\begin{lstlisting}[language=Python]
hmm_model = GaussianHMM(n_components=3, covariance_type="full", n_iter=200, random_state=42)
hmm_model.fit(returns)
hidden_states = hmm_model.predict(returns)                 # Viterbi
state_order = np.argsort([np.sqrt(hmm_model.covars_[i][0][0]) for i in range(3)])
\end{lstlisting}

\section{7. Results}\label{results}

{\def\LTcaptype{none} % do not increment counter
\begin{longtable}[]{@{}
  >{\raggedright\arraybackslash}p{(\linewidth - 2\tabcolsep) * \real{0.5000}}
  >{\raggedright\arraybackslash}p{(\linewidth - 2\tabcolsep) * \real{0.5000}}@{}}
\toprule\noalign{}
\begin{minipage}[b]{\linewidth}\raggedright
Item
\end{minipage} & \begin{minipage}[b]{\linewidth}\raggedright
Value
\end{minipage} \\
\midrule\noalign{}
\endhead
\bottomrule\noalign{}
\endlastfoot
EM converged & True \\
Model log-likelihood & 6019.79 (1,859 days) \\
Regimes & Bull (low vol), Sideways (medium), Bear (high vol) after
sorting \\
\end{longtable}
}

The colored DAX chart shows the model tracking real episodes: the 1990s
bull phases, corrections such as 1994, and the high-volatility cluster
around the 1998 crisis. The transition-matrix diagonal is dominant,
confirming regime persistence.

\section{8. Discussion and Conclusion}\label{discussion-and-conclusion}

The 3-state Gaussian HMM extracts interpretable, persistent market
regimes from real data without labels. Its main assumptions are the
Markov property (next state depends only on the current one) and
Gaussian emissions; return fat tails and regime duration dependence
motivate extensions (t-emissions, semi-Markov models).

\section{9. Answers to Viva Questions}\label{answers-to-viva-questions}

\begin{itemize}
\tightlist
\item
  \textbf{What is hidden in an HMM?} The state sequence: we observe
  emissions (returns) but not the underlying regime that generated them.
\item
  \textbf{Viterbi vs forward algorithm?} Viterbi returns the single most
  likely state path (max over paths); the forward algorithm computes the
  total probability P(X \textbar{} λ) by summing over all paths.
\item
  \textbf{What is the Markov assumption?} The future state depends on
  the past only through the present state: P(z\_t \textbar{}
  z\_\{1:t−1\}) = P(z\_t \textbar{} z\_\{t−1\}).
\end{itemize}

\begin{center}\rule{0.5\linewidth}{0.5pt}\end{center}

\section{10. Complete Code Listing}\label{complete-code-listing}

\emph{Full runnable program for this experiment, extracted from
\passthrough{\lstinline!notebooks/08\_hidden\_markov\_model.ipynb!}
(executed; results above are its actual output).}

\textbf{Listing 16.1 --- Core libraries}

\begin{lstlisting}[language=Python]
# ---- Core libraries ----
import numpy as np
import pandas as pd
import matplotlib.pyplot as plt
import seaborn as sns
from hmmlearn.hmm import GaussianHMM

# Styling setup
plt.style.use('seaborn-v0_8-whitegrid' if 'seaborn-v0_8-whitegrid' in plt.style.available else 'default')
plt.rcParams['font.sans-serif'] = 'DejaVu Sans'
plt.rcParams['figure.figsize'] = (10, 6)
plt.rcParams['figure.dpi'] = 110
np.random.seed(42)

# ---- Load the real DAX stock index (daily closes 1991-1998) ----
df_market = pd.read_csv('../data/stock_index.csv')
dates = pd.to_datetime(df_market['Date'])
print(f"Trading days: {len(df_market)} | {dates.iloc[0].date()} to {dates.iloc[-1].date()}")
print(df_market.head(3).to_string(index=False))

returns = df_market['Daily_Return'].values.reshape(-1, 1)   # HMM observations
prices = df_market['Close'].values                          # only used for visualization
\end{lstlisting}

\textbf{Listing 16.2 --- Fit a 3-state Gaussian HMM with the Baum-Welch
(EM) algorithm}

\begin{lstlisting}[language=Python]
# ---- Fit a 3-state Gaussian HMM with the Baum-Welch (EM) algorithm ----
n_components = 3
hmm_model = GaussianHMM(
    n_components=n_components,        # Bull / Sideways / Bear regimes
    covariance_type="full",           # each state has its own full covariance Gaussian emission
    n_iter=200,                       # EM iterations
    random_state=42,
    verbose=False
)

hmm_model.fit(returns)
print("EM Training Converged:", hmm_model.monitor_.converged)
print(f"Model Log-Likelihood: {hmm_model.score(returns):.2f}")

# ---- Decode the most likely hidden state sequence with the Viterbi algorithm ----
hidden_states = hmm_model.predict(returns)

# ---- Sort regimes by volatility so labels have a consistent economic meaning ----
# Regime 0: Low Volatility (Bull) | Regime 1: Moderate (Sideways) | Regime 2: High (Bear/Crisis)
volatilities = [np.sqrt(hmm_model.covars_[i][0][0]) for i in range(n_components)]
state_order = np.argsort(volatilities)                                           # ascending volatility
state_mapping = {old_state: new_state for new_state, old_state in enumerate(state_order)}

sorted_hidden_states = np.array([state_mapping[s] for s in hidden_states])       # relabel sequence
sorted_means = [hmm_model.means_[old][0] for old in state_order]                 # reordered means
sorted_stds  = [np.sqrt(hmm_model.covars_[old][0][0]) for old in state_order]    # reordered std devs

# ---- Reorder the transition matrix to match the sorted regimes ----
sorted_transmat = np.zeros((n_components, n_components))
for i in range(n_components):
    for j in range(n_components):
        sorted_transmat[i, j] = hmm_model.transmat_[state_order[i], state_order[j]]

regime_names = ["Bull (Low Vol)", "Sideways (Med Vol)", "Bear (High Vol)"]
\end{lstlisting}

\textbf{Listing 16.3 --- Transition probability matrix}

\begin{lstlisting}[language=Python]
# ---- Transition probability matrix ----
fig, ax = plt.subplots(figsize=(6, 4.5))
sns.heatmap(sorted_transmat, annot=True, fmt='.3f', cmap='Blues', ax=ax,
            xticklabels=regime_names, yticklabels=regime_names, cbar=False)
ax.set_title("HMM Transition Probability Matrix ($A_{ij}$)", fontweight='bold')
ax.set_xlabel("Transition To Regime")
ax.set_ylabel("Current Regime")
plt.tight_layout()
plt.show()

# ---- Viterbi-decoded regimes over the price series ----
colors = ['forestgreen', 'royalblue', 'crimson']
x = np.arange(len(prices))

fig, (ax1, ax2) = plt.subplots(2, 1, figsize=(13, 7), sharex=True, gridspec_kw={'height_ratios': [3, 1]})

# Top: color every consecutive price segment by its decoded regime
for i in range(len(prices) - 1):
    ax1.plot([x[i], x[i + 1]], [prices[i], prices[i + 1]], color=colors[sorted_hidden_states[i]], lw=1.2)
for i, name in enumerate(regime_names):
    ax1.plot([], [], color=colors[i], label=name, lw=2.5)
ax1.set_title("DAX Index 1991-1998 Colored by Decoded HMM Regime (Viterbi Path)", fontsize=13, fontweight='bold')
ax1.set_ylabel("Index Level")
ax1.legend(loc='upper left')

# Bottom: raw sequence of decoded states over time
ax2.scatter(x, sorted_hidden_states, c=[colors[s] for s in sorted_hidden_states], s=8, marker='|')
ax2.set_yticks(range(n_components))
ax2.set_yticklabels(regime_names)
ax2.set_xlabel("Trading Day Index")

# Label the x-axis with years (real dates)
year_pos = np.where(dates.dt.year.diff().fillna(0).values != 0)[0]
ax2.set_xticks(year_pos)
ax2.set_xticklabels(dates.dt.year.iloc[year_pos].astype(str))

plt.tight_layout()
plt.show()
\end{lstlisting}
